\section*{Chapter \ref{ch:b3:topology} --- General Topology}

\begin{solution}{exo:b3:topology:1}
(a) Four topologies on $\{a,b\}$: indiscrete
$\{\varnothing, X\}$; discrete; the two \emph{Sierpi\'nski}
topologies $\{\varnothing, \{a\}, X\}$ and $\{\varnothing,
\{b\}, X\}$. The last two are homeomorphic (swap $a, b$): three
classes. Connected: all except the discrete one (only the
discrete topology contains a proper nonempty clopen set).
Hausdorff: only the discrete one (in the others, $b$'s only
neighborhood is $X$, or $a$'s).

(b) $\Q$: interior $\varnothing$ (every interval contains
irrationals), closure $\R$ (density), boundary $\R$. $[0,1)
\cup \{2\}$: interior $(0,1)$, closure $[0,1]\cup\{2\}$,
boundary $\{0, 1, 2\}$. $\{1/n\}$: interior $\varnothing$,
closure $\{1/n\}\cup\{0\}$, boundary the closure itself.
\end{solution}

\begin{solution}{exo:b3:topology:2}
(a) Every open $V \subseteq Y$ is a union $\bigcup B_j$ of basis
sets, and $f^{-1}(V) = \bigcup f^{-1}(B_j)$: if the latter are
open, so is the former; the converse is trivial.

(b) $f^{-1}\bigl(\intoo{\frac12}{\frac32}\bigr) = [0, \infty)$
is not open: $f$ is not continuous; right-continuity at each
point is clear ($f$ is locally constant on the right of every
point). The sets $[a, b)$ satisfy the basis criterion
(\cref{def:b3:topology:basis}): $[a,b)\cap[c,d) = [\max(a,c),
\min(b,d))$. The Sorgenfrey topology is finer than the usual
one, since $(a, b) = \bigcup_{a < c < b}[c, b)$; strictly:
$[0,1)$ is Sorgenfrey-open, not usual-open. Continuity from the
Sorgenfrey line at $x$: a Sorgenfrey neighborhood of $x$
contains a basis set $[a, b) \ni x$, hence contains $[x, b)$ ---
so the continuity condition reads: for every $\varepsilon > 0$
there is $\delta > 0$ with $\abs{f(t) - f(x)} < \varepsilon$
whenever $x \leq t < x + \delta$. That is exactly
right-continuity at $x$.
\end{solution}

\begin{solution}{exo:b3:topology:3}
Two nonempty opens have finite complements, so their
intersection has finite complement: nonempty (the set is
infinite) --- no two points have disjoint neighborhoods: not
Hausdorff, hence not metrizable
(\cref{def:b3:topology:hausdorff}). Let $(x_n)$ be injective and
$x$ arbitrary: a neighborhood of $x$ contains a cofinite open
$U$; the finitely many points of $X\setminus U$ are hit by at
most finitely many indices (injectivity), so a tail of the
sequence lies in $U$: $x_n \to x$, for \emph{every} $x$. The
uniqueness proof needs two disjoint neighborhoods to separate
two alleged limits --- precisely what fails here.
\end{solution}

\begin{solution}{exo:b3:topology:4}
(a) Continuity is by construction
(\cref{def:b3:topology:constructions}). Openness: an open $W
\subseteq X \times Y$ is a union of boxes $U_i\times V_i$, and
$\pi_X(W) = \bigcup U_i$ (nonempty boxes project onto their
factors), open. Not closed: $H = \{(x,y) : xy = 1\}$ is closed
in $\R^2$ (preimage of $\{1\}$ under the continuous product),
but $\pi_X(H) = \R\setminus\{0\}$ is not closed.

(b) ($\Rightarrow$) Projections are continuous. ($\Leftarrow$)
Let $x^{(k)} \to x$ componentwise and $W = \prod U_n$ a basis
neighborhood of $x$, with $U_n = X_n$ except for $n \in F$
finite. For each $n \in F$, $x_n^{(k)} \in U_n$ for $k \geq
K_n$; for $k \geq \max_F K_n$, $x^{(k)} \in W$. The formula $d$
defines a metric (each summand is one, up to the standard
verification that $\min(1, d_n)$ is); its balls: $B_d(x,
2^{-N})$ contains the basis box $\{y : d_n(x_n,y_n) < 2^{-N},\
n \leq N'\}$ for suitable $N'$ (the tail $\sum_{n > N'} 2^{-n}$
is small), and conversely every basis box contains a $d$-ball:
the two topologies have the same neighborhoods of each point.

(c) In the product topology, $x^{(k)} \to 0$ by (b). The box-open
set $U = \prod_n \intoo{-\tfrac1n}{\tfrac1n}$ contains $0$, but
$x^{(k)} \notin U$ for every $k$ (the $n$-th coordinate $1/k
\geq 1/n$ when $n \geq k$): no tail enters $U$. The box topology
is strictly finer and not a product-type topology for
convergence purposes.
\end{solution}

\begin{solution}{exo:b3:topology:5}
$\varphi \colon \R \to S^1$, $x \mapsto \eu^{2\iu\pi x}$, is
continuous, surjective, and constant on classes mod $\Z$: it
induces a continuous bijection $\bar\varphi \colon \R/\Z \to
S^1$ (\cref{prop:b3:topology:universal}(b)). The projection
$\pi$ is \emph{open}: for $U \subseteq \R$ open,
$\pi^{-1}(\pi(U)) = \bigcup_{n}(U + n)$ is open, so $\pi(U)$ is
open. Hausdorff: let $\bar x \neq \bar y$; choose
representatives with $\delta = \min_{n}\abs{x - y - n} > 0$;
the images of the intervals of radius $\delta/3$ around $x$ and
$y$ are open (openness of $\pi$), contain $\bar x, \bar y$, and
are disjoint (two preimage points would be $< 2\delta/3$ apart
mod $\Z$). Compact: $\R/\Z = \pi([0,1])$, a continuous image of
a compact space (\cref{thm:b3:topology:compactprops}(2)). Now
$\bar\varphi$ is a continuous bijection from a compact space to
the Hausdorff $S^1$: a homeomorphism
(\cref{cor:b3:topology:compacthomeo}).

For (iii): the composite $[0,1] \hookrightarrow \R
\xrightarrow{\pi} \R/\Z$ is continuous, surjective, and
identifies exactly $0 \sim 1$: it induces a continuous bijection
$[0,1]/(0\sim1) \to \R/\Z$ from a compact space (continuous
image of $[0,1]$ under the quotient projection) to a Hausdorff
one: a homeomorphism. Composing: all three spaces are
homeomorphic.
\end{solution}

\begin{solution}{exo:b3:topology:6}
(a) A cover of $K_1 \cup \dots \cup K_m$ restricts to a cover of
each $K_j$: finitely many opens per piece suffice. An
intersection $\bigcap K_i$ is closed in the compact $K_{i_0}$
(compacts are closed in the ambient metric space), hence
compact.

(b) $x \mapsto d(x, F)$ is continuous ($1$-Lipschitz) and
strictly positive on $K$ ($d(x, F) = 0$ means $x \in \bar F =
F$); on the compact $K$ it attains a minimum $m > 0$: $d(K, F)
\geq m$. Counterexample without compactness: $F_1 = \{n : n \geq
2\}$ and $F_2 = \{n + \frac1n : n \geq 2\}$ are closed,
disjoint, at distance $\inf_n \frac1n = 0$.

(c) If $X$ is compact, every continuous $f$ is bounded
(\cref{thm:b3:topology:compactprops}(2)). Conversely, if $X$ is
not compact, take $(x_n)$ with no convergent subsequence
(\cref{thm:b3:topology:metriccompact}); passing to a
subsequence we may assume the $x_n$ pairwise distinct, and no
point of $X$ is a cluster point of the sequence, so
\[
r_n = \min\Bigl(2^{-n},\ \tfrac13\,d\bigl(x_n, \{x_m : m \neq
n\}\bigr)\Bigr) > 0,
\]
and the balls $B(x_n, r_n)$ are pairwise disjoint (a common
point of the $n$-th and $m$-th would give $d(x_n, x_m) < r_n +
r_m \leq \frac23 d(x_n, x_m)$). Define
\[
f(x) = \sum_{n} n\,\max\Bigl(0,\ 1 - \frac{d(x,
x_n)}{r_n}\Bigr):
\]
at each $x$ at most one summand is nonzero, and $f(x_n) = n$.
Continuity at $x$: let $y_k \to x$; each value $f(y_k)$ is
either $0$ or a single bump value $\mathrm{bump}_{n_k}(y_k)$.
If some index $n$ occurs infinitely often, along that
subsequence $f(y_k) = \mathrm{bump}_n(y_k) \to
\mathrm{bump}_n(x)$, which equals $f(x)$ (if $x \in B(x_n,
r_n)$, the whole tail lies in this open ball and no other
index occurs; if $x \notin B(x_n, r_n)$ then
$\mathrm{bump}_n(x) = 0 = f(x)$, since $x$, adherent to the
$n$-th ball, lies in no other open ball). If $n_k \to \infty$:
$d(y_k, x_{n_k}) < r_{n_k} \leq 2^{-n_k} \to 0$, so $x_{n_k}
\to x$, making $x$ a cluster point of $(x_n)$ --- excluded; so
this case concerns finitely many $k$ only. In every case
$f(y_k) \to f(x)$: $f$ is continuous, and unbounded.
\end{solution}

\begin{solution}{exo:b3:topology:7}
Suppose no $\delta$ works: for each $n$ there is $A_n$ of
diameter $< 1/n$ contained in no single $U_i$; pick $x_n \in
A_n$. By compactness (\cref{thm:b3:topology:metriccompact}),
a subsequence $x_{n_k} \to x$; pick $i$ and $r > 0$ with $B(x,
r) \subseteq U_i$. For large $k$: $d(x_{n_k}, x) < r/2$ and
$1/n_k < r/2$, so $A_{n_k} \subseteq B(x_{n_k}, 1/n_k)
\subseteq B(x, r) \subseteq U_i$ --- contradiction. Heine:
given $\varepsilon$, cover $Y$ by balls $B(y, \varepsilon/2)$;
the preimages form an open cover of $X$; let $\delta$ be a
Lebesgue number: if $d(x, x') < \delta$, the pair $\{x, x'\}$
has diameter $< \delta$, lies in one preimage, and $d(f(x),
f(x')) < \varepsilon$.
\end{solution}

\begin{solution}{exo:b3:topology:8}
(a) $GL_n(\R) = \det^{-1}(\R\setminus\{0\})$ is open ($\det$ is
polynomial, continuous). Disconnected: $\det$ maps it onto
$\R\setminus\{0\}$, and a connected space has connected
continuous images (\cref{thm:b3:topology:connectedbasics}(2));
$\R\setminus\{0\}$ is not an interval. $GL_n(\C)$: for $A, B$
invertible, $p(\lambda) = \det\bigl((1-\lambda)A + \lambda
B\bigr)$ is a polynomial in $\lambda$ with $p(0), p(1) \neq 0$,
hence not identically zero: it has finitely many roots in $\C$.
The plane minus finitely many points is path-connected (avoid
the points by going around), so there is a path $\lambda(t)$
from $0$ to $1$ with $p(\lambda(t)) \neq 0$: $t \mapsto
(1-\lambda(t))A + \lambda(t)B$ is a path in $GL_n(\C)$.

(b) $\det(A + \varepsilon I) = \chi_A(-\varepsilon)\cdot(\pm1)$
vanishes for at most $n$ values of $\varepsilon$: invertible
matrices $A + \varepsilon I \to A$ as $\varepsilon \to 0$:
density.
\end{solution}

\begin{solution}{exo:b3:topology:9}
(a) Let $A \subseteq \Q$ contain $p < q$ and pick an irrational
$\alpha \in (p, q)$: $A = (A \cap (-\infty, \alpha)) \sqcup
(A\cap(\alpha, +\infty))$ splits $A$ into two nonempty relatively
open pieces: $A$ is disconnected. Components are singletons.
Local compactness fails: a compact neighborhood $V$ of $0$ in
$\Q$ contains $[-r, r]\cap\Q$ for some $r > 0$, which is closed
in $V$, hence compact; but a sequence of rationals in $[-r,r]$
converging (in $\R$) to an irrational has no subsequence
converging \emph{in $\Q$}: contradiction with
\cref{thm:b3:topology:metriccompact}.

(b) Let $\gamma = (\gamma_1, \gamma_2) \colon [0,1] \to S$ be a
path with $\gamma(0) = (0,0)$ and $\gamma(1) \in \Gamma$. The
set $\{t : \gamma_1(t) = 0\}$ is closed and does not contain
$1$; let $t^*$ be its supremum, so $\gamma_1(t^*) = 0$ and
$\gamma_1 > 0$ on $(t^*, 1]$. By continuity of $\gamma_2$ at
$t^*$, choose $\delta$ with $\abs{\gamma_2(t) -
\gamma_2(t^*)} < \tfrac12$ for $t \in [t^*, t^* + \delta]$. Fix
$t_0 = t^* + \delta$ (may assume $\leq 1$): $\gamma_1(t_0) > 0$,
and $\gamma_1$ takes every value of $(0, \gamma_1(t_0)]$ on
$[t^*, t_0]$ (intermediate value theorem,
\cref{thm:b3:topology:connectedbasics}). Choose $k$ large with
$a_k = \frac1{\pi/2 + 2k\pi} < \gamma_1(t_0)$ and $b_k =
\frac1{3\pi/2 + 2k\pi} < \gamma_1(t_0)$: there are $s, s' \in
(t^*, t_0]$ with $\gamma_1(s) = a_k$, $\gamma_1(s') = b_k$;
since the points lie on the graph, $\gamma_2(s) = \sin(1/a_k) =
1$ and $\gamma_2(s') = -1$. Both cannot be within $\frac12$ of
$\gamma_2(t^*)$: contradiction. $S$ is connected but not
path-connected.
\end{solution}

\begin{solution}{exo:b3:topology:10}
(a) $C_n$ consists exactly of the $x$ admitting a ternary
expansion with digits in $\{0, 2\}$ up to rank $n$ (induction:
removing middle thirds removes first digit $1$, etc.; endpoints
have two expansions, one avoiding $1$s), so $C = \bigcap C_n$
is the set of sums $\sum a_n3^{-n}$, $a_n \in \{0,2\}$. Compact:
each $C_n$ is a finite union of closed intervals; $C$ is closed
in $[0,1]$. Empty interior: $C \subseteq C_n$, a union of
intervals of length $3^{-n}$; an interior interval of length
$\ell > 0$ would fit in one of them for all $n$. Perfect: given
$x \in C$ and $m$, flip the digit $a_m$: the new point is in
$C$, distinct, within $2\cdot3^{-m}$ of $x$.

(b) If $x \neq y \in C$, their expansions first differ at some
rank $k$; between them lies a removed middle-third interval
(the gap at rank $k$ separating digit $0$ from digit $2$),
providing a point $c \notin C$, $x < c < y$ (say): $C =
(C \cap (-\infty, c)) \sqcup (C \cap (c, \infty))$ splits any
subset containing both points. Components are singletons.

(c) The digit map $\beta\colon x \mapsto (a_n(x))$ is a
bijection onto $\{0,2\}^\N$ (existence and uniqueness of
$\{0,2\}$-expansions: distinct digit sequences give points at
distance $\geq 3^{-k} > 0$ where they first differ, as in
\cref{pb:b3:topology:1}, question 1). It is continuous:
$\abs{x - y} < 3^{-m}$ forces agreement of the first $m$ digits
(same computation), so $\beta$ maps small balls into basis
boxes. A continuous bijection from the compact $C$ to the
Hausdorff product is a homeomorphism
(\cref{cor:b3:topology:compacthomeo}; the product is Hausdorff:
separate at a differing coordinate). Uncountability: Cantor's
diagonal on $\{0,2\}^\N$. Finally $\{0,2\}^\N \times
\{0,2\}^\N \cong \{0,2\}^\N$ by interleaving digits (a
homeomorphism: componentwise continuity both ways), so $C
\times C \cong C$.
\end{solution}

\begin{solution}{exo:b3:topology:11}
For $y \in \R^n$, set
\[
\tau(y) = \Bigl(\frac{2y}{\norm y^2 + 1},\ \frac{\norm y^2 -
1}{\norm y^2 + 1}\Bigr) \in \R^n\times\R = \R^{n+1}.
\]
One checks $\norm{\tau(y)} = 1$, $\tau(y) \neq N$, and
$\sigma(\tau(y)) = y$, $\tau(\sigma(x)) = x$: $\sigma$ and
$\tau$ are mutually inverse, both continuous (rational formulas
with nonvanishing denominators): $S^n\setminus\{N\} \cong
\R^n$. Extend to $\hat\sigma \colon S^n \to \widehat{\R^n}$ by
$N \mapsto \infty$: a bijection, continuous at every point of
$S^n\setminus\{N\}$, and at $N$: a basis neighborhood of
$\infty$ is $\widehat{\R^n}\setminus K$, $K$ compact, $K
\subseteq \bar B(0, R)$; since $\norm{\sigma(x)}^2 = \frac{1 +
x_{n+1}}{1 - x_{n+1}} \to \infty$ as $x_{n+1} \to 1$, the set
$\{x \in S^n : x_{n+1} > 1 - \eta\}$ maps outside $\bar B(0,R)$
for small $\eta$: continuity. A continuous bijection from the
compact $S^n$ to the Hausdorff $\widehat{\R^n}$ is a
homeomorphism. Removing another point $p$: a rotation of $S^n$
maps $p$ to $N$ (rotations are homeomorphisms), reducing to the
computed case: $S^n \setminus \{p\} \cong \R^n$.
\end{solution}

\begin{solution}{pb:b3:topology:1}
\textbf{1.} Suppose the expansions of $x, y \in C$ first differ
at rank $k \leq m$, say $b_k(x) = 1$, $b_k(y) = 0$. Then
\[
x - y \;\geq\; \frac{2}{3^k} - \sum_{j > k}\frac{2}{3^j}
= \frac{2}{3^k} - \frac{1}{3^k} = \frac1{3^k} \geq \frac1{3^m}:
\]
contrapositive: $\abs{x - y} < 3^{-m}$ forces agreement up to
rank $m$.

\textbf{2.} By question 1, $b_n$ is constant on $C \cap (x -
3^{-n}, x + 3^{-n})$: locally constant, hence continuous. The
map $x \mapsto (b_n(x))_n \in \{0,1\}^\N$ is continuous
(componentwise, \cref{prop:b3:topology:universal}(a)) and
bijective (unique $\{0,2\}$-digit expansions); from the compact
$C$ to a Hausdorff space, it is a homeomorphism
(\cref{cor:b3:topology:compacthomeo}).

\textbf{3.} Continuity: if $\abs{x - y} < 3^{-m}$, the first
$m$ digits agree, so $\abs{g(x) - g(y)} \leq \sum_{n > m}2^{-n}
= 2^{-m}$. Surjectivity: every $t \in [0,1]$ has a binary
expansion $t = \sum b_n2^{-n}$, and $t = g\bigl(\sum
2b_n3^{-n}\bigr)$. Not injective: $g$ identifies the two Cantor
points with digits $(0,1,1,1,\dots)$ and $(1,0,0,\dots)$ ---
both map to $\frac12$ (the dyadic ambiguity $0.0111\ldots_2 =
0.1000\ldots_2$).

\textbf{4.} Each component of $\Phi$ is continuous by the same
estimate (its digits are a subsequence of the $b_n$).
Surjectivity: given $(s, t) \in [0,1]^2$, choose binary digits
$(c_n)$ of $s$ and $(d_n)$ of $t$, and interleave: the point $x
\in C$ with $b_{2n-1}(x) = c_n$, $b_{2n}(x) = d_n$ satisfies
$\Phi(x) = (s,t)$.

\textbf{5.} The gaps are pairwise disjoint with endpoints in
$C$, so the formula defines $f$ unambiguously on $[0,1]$; at
the endpoints of a gap the affine formula returns $\Phi(u)$,
$\Phi(v)$: consistent with $f = \Phi$ on $C$. Surjectivity:
already $f(C) = \Phi(C) = [0,1]^2$.

\textbf{6.} At $t_0 \notin C$: $t_0$ lies in an open gap on
which $f$ is affine: continuous. At $x \in C$: we first record
two estimates.

\emph{(i) On $C$}: if $\abs{y - x} \leq 3^{-M}$, $y \in C$,
then the digits agree up to $M$, so each component of $\Phi(y)
- \Phi(x)$ is at most $\sum_{n > \lfloor M/2\rfloor} 2^{-n} =
2^{-\lfloor M/2\rfloor}$.

\emph{(ii) Across a gap}: a gap $(u,v)$ removed at stage $k$
has length $3^{-k}$, and its endpoints have digits agreeing up
to rank $k - 1$ (they differ from rank $k$ on), so
$\norm{\Phi(v) - \Phi(u)}_\infty \leq
2^{-\lfloor (k-1)/2\rfloor}$.

Now let $\abs{t - x} < 3^{-M}$ with $t$ in a gap $(u, v)$ of
stage $k$, and say $u$ is the endpoint on $x$'s side, so
$\abs{u - x} < 3^{-M}$ and $0 < t - u < 3^{-M}$. Then
\[
\norm{f(t) - \Phi(x)}_\infty
\leq \norm{\Phi(u) - \Phi(x)}_\infty
+ \frac{t - u}{v - u}\,\norm{\Phi(v) - \Phi(u)}_\infty
\leq 2^{-\lfloor M/2\rfloor} + \min\bigl(1,\ 3^{k - M}\bigr)\,
2^{-\lfloor(k-1)/2\rfloor} .
\]
For $k \geq M$ the last term is $\leq 2^{-\lfloor(M-1)/2
\rfloor}$; for $k < M$, since $2^{-\lfloor(k-1)/2\rfloor} \leq
\sqrt2\,\cdot 2^{-k/2}$,
\[
3^{k-M}\,2^{-\lfloor(k-1)/2\rfloor}
\leq \sqrt2\,\Bigl(\frac{3}{\sqrt2}\Bigr)^{k} 3^{-M}
\leq \sqrt2\,\Bigl(\frac{3}{\sqrt2}\Bigr)^{M} 3^{-M}
= \sqrt2\; 2^{-M/2}
\]
(the middle quantity increases in $k$). In all cases
$\norm{f(t) - \Phi(x)}_\infty \leq C\,2^{-M/2}$
with an absolute constant $C$: letting $M \to \infty$ proves
continuity at $x$ (points $t \in C$ are covered by (i)). Hence
$f$ is a continuous surjection $[0,1] \to [0,1]^2$ --- indeed
the estimates show it is H\"older of exponent $\log 2/(2\log
3)$-flavor, but continuity is all we claimed.

\textbf{7.} For $[0,1]^n$: split the digits of $x \in C$ into
$n$ interleaved subsequences and repeat questions 4--6
verbatim. For $\R \to \R^2$: let $h \colon [0,1] \to [-1,1]^2$
be a continuous surjection (rescale $f$). Define $F$ on $[m, m
+ \frac12]$ ($m \in \Z$) as a copy of $h$ rescaled to fill
$[-(m+1), m+1]^2$ for $m \geq 0$ (and symmetrically for $m <
0$), and on $[m + \frac12, m + 1]$ as the affine segment
joining the endpoint values: $F$ is continuous (gluing on
closed pieces, \cref{prop:b3:topology:gluing}(b)) and its image
contains $\bigcup_m [-(m+1), m+1]^2 = \R^2$.

\textbf{8.} $[0,1]$ is compact and $[0,1]^2$ Hausdorff: a
continuous injection is a homeomorphism onto its image
(\cref{cor:b3:topology:compacthomeo} applied to the corestriction).

\textbf{9.} Remove $\frac12$: $[0,1]\setminus\{\frac12\}$ is
disconnected. If $h \colon [0,1] \to [0,1]^2$ were a
homeomorphism, $[0,1]^2\setminus\{h(\frac12)\}$ would be
disconnected too (homeomorphic images of disconnected spaces
are disconnected); but the square minus a point is
path-connected: join two points by a two-segment path avoiding
the puncture. Contradiction: $[0,1] \not\cong [0,1]^2$.

\textbf{10.} By question 8, an injective continuous surjection
$[0,1] \to [0,1]^2$ would be a homeomorphism, contradicting
question 9. Compactness entered exactly in
\cref{cor:b3:topology:compacthomeo}: it is what makes the
continuous bijection's inverse continuous (images of closed
sets are compact, hence closed). Without compactness the
conclusion genuinely fails: $[0, 2\pi) \to S^1$ is a continuous
bijection that is not a homeomorphism.

\textbf{11.} Proved here: \emph{there is a continuous
surjection, but no continuous bijection, from $[0,1]$ onto
$[0,1]^2$; in particular $[0,1] \not\cong [0,1]^2$.} So
``dimension'' is not preserved by continuous surjections ---
cardinality and even continuity cannot see it --- but it
\emph{is} a topological invariant at the level of
homeomorphisms, at least for $1$ versus $2$ dimensions, where
connectedness-after-deletion suffices. The general statement
--- \emph{invariance of domain}: $\R^m \cong \R^n$ implies $m =
n$, and a continuous injection $\R^n \to \R^n$ is open ---
is true but requires algebraic topology (homology), beyond
this course.
\end{solution}
